% !TEX root = ../main.tex
\section{Three Fundamental Estimators and Similarity}
\label{sec:algorithms-and-errors}

\subsection{SGD, MSGD, and TSGD}
\label{subsec:algorithms-and-assumptions}

After taking a step to $x_{k+1}$, we need a gradient estimate to
determine the next update. The identity
\[
\nabla f(x_{k+1})
=\nabla f(x_k)+\bigl[\nabla f(x_{k+1})-\nabla f(x_k)\bigr]
\]
suggests two approaches. SGD estimates the left-hand side directly
by drawing a fresh stochastic gradient at the new point. TSGD estimates
the right-hand side: its current estimate $g_k$ already approximates
the first term, so it needs new information only for the change.
Using a stochastic gradient oracle $h$ with the same fresh sample
at both endpoints gives
\[
g_{k+1}
=\underbrace{g_k}_{\approx\,\nabla f(x_k)}
+\underbrace{h(x_{k+1},\xi_{k+1})-h(x_k,\xi_{k+1})}
 _{\approx\,\nabla f(x_{k+1})-\nabla f(x_k)}.
\]
Thus TSGD updates the gradient information it already has.
The sample $\xi_{k+1}$ is shared within the difference and drawn afresh
at each iteration. For a finite sum, this means evaluating the same
randomly selected component at both points. TSGD retains the information
in $g_k$, including its estimation error; the effect of these inherited
errors is examined in Section~\ref{subsec:martingale-interpretation}.

We formalize the setup as follows.
\begin{assumption}[Smooth, lower-bounded objective]
\label{ass:smooth-objective}
The function $f:\mathbb R^d\to\mathbb R$ is lower bounded,
$f_{\inf}:=\inf_x f(x)>-\infty$, and has an $L$-Lipschitz gradient
with $L>0$:
\[
\|\nabla f(x)-\nabla f(y)\|\leq L\|x-y\|
\qquad\text{for all }x,y\in\mathbb R^d.
\]
\end{assumption}
\begin{assumption}[Unbiased stochastic gradients]
\label{ass:unbiased-oracle}
The oracle satisfies $\mathbb E_{\xi}[h(x,\xi)]=\nabla f(x)$ for every $x$.
Each new sample is independent of the preceding history.
\end{assumption}

All three methods use positive deterministic stepsizes $\eta_k$ and
the same iterate update, starting from a deterministic $x_0$:
\begin{equation}
\label{eq:tsgd-iterate}
x_{k+1}=x_k-\eta_k g_k.
\end{equation}
Their gradient estimates differ:
\begin{align}
\mathrm{SGD}:\quad g_{k+1}
  &=h(x_{k+1},\xi_{k+1}), \label{eq:sgd-estimate}\\
\mathrm{MSGD}:\quad g_{k+1}
  &=(1-\beta_k)g_k+\beta_k h(x_{k+1},\xi_{k+1}),
  \quad 0<\beta_k\leq1, \label{eq:msgd-estimate}\\
\mathrm{TSGD}:\quad g_{k+1}
  &=g_k+D_{k+1}. \label{eq:tsgd-estimate}
\end{align}
In the standard pointwise-oracle construction,
\begin{equation}
\label{eq:shared-sample-increment}
D_{k+1}=h(x_{k+1},\xi_{k+1})-h(x_k,\xi_{k+1}).
\end{equation}
SGD replaces its estimate with a fresh stochastic gradient at every
iteration. MSGD averages fresh and past estimates to reduce noise, but
can lag behind a changing gradient. TSGD accounts for the gradient change
through transport.
Unless specified otherwise, SGD and MSGD initialize with
$g_0=h(x_0,\xi_0)$. TSGD permits any $g_0$ with finite mean-square
error, including a biased estimate.

\subsection{Understanding Similarity}
\label{subsec:understanding-similarity}

SGD is typically analyzed under a bounded-variance assumption.
Write $r(x,\xi)=h(x,\xi)-\nabla f(x)$ for the gradient noise.
\begin{assumption}[Bounded variance (BV)]
\label{ass:bounded-variance}
For some $\sigma\geq0$,
\begin{equation}
\label{eq:bounded-variance}
\mathbb E_{\xi}[\|r(x,\xi)\|^2]\leq\sigma^2
\qquad\text{for all }x.
\end{equation}
\end{assumption}
Transport instead uses a condition on gradient differences.
\begin{assumption}[Centered $\delta$-similarity]
\label{ass:oracle-similarity}
For some $\delta\geq0$,
\begin{equation}
\label{eq:oracle-similarity}
\mathbb E_{\xi}[\|r(x,\xi)-r(y,\xi)\|^2]
\leq\delta^2\|x-y\|^2
\qquad\text{for all }x,y.
\end{equation}
\end{assumption}
Under uniform finite-sum sampling, this condition coincides with the
Hessian-variance condition of
\citet{szlendak2021permutationcompressorsprovablyfaster} and the
second-order similarity assumption of
\citet{khaled2023fasterfederatedoptimizationsecondorder}.
More recently, \citet{pmlr-v306-shulgin26b} generalized it through
$\gamma$-similarity, which recovers this condition at $\gamma=1$.
We use the name \emph{centered $\delta$-similarity} to emphasize that it
measures the variation of gradient noise rather than its magnitude.

Here $\sigma$ bounds the magnitude of individual gradient errors, whereas
$\delta$ measures their variation under shared sampling. A small $\delta$
means that stochastic gradient differences closely match deterministic
gradient differences; it does not require the gradients themselves, or
their changes, to be small. Similarity is a condition on the stochastic
oracle and applies to both finite-sum and population objectives.

For a finite sum $f=\frac1n\sum_{i=1}^n f_i$ with quadratic components
$f_i(x)=\tfrac12x^\top A_i x+b_i^\top x+c_i$ and symmetric $A_i$, the linear terms
$b_i^\top x$ cancel from gradient differences, so identical Hessians
give exact transport even when individual gradients differ greatly.
Hessian variation controls $\delta$, whereas $L$ concerns the average
Hessian; neither must be smaller than the other.
Appendix~\ref{app:quadratic-similarity} gives the calculation.

Scalar least squares separates the two quantities. Let
$f(x)=\tfrac12\mathbb E[(ux-b)^2]$, with feature $u$ satisfying
$0<\mathbb E[u^2]$ and $\mathbb E[u^4]<\infty$, and label
$b=ux^\star+\varepsilon$. The label noise $\varepsilon$ is independent
of $u$, with $\mathbb E[\varepsilon]=0$ and
$\mathbb E[\varepsilon^2]=s^2<\infty$. The stochastic gradient is
$h(x;u,\varepsilon)=u^2(x-x^\star)-u\varepsilon$, so
\begin{equation}
\label{eq:least-squares-similarity}
\begin{aligned}
h(x;u,\varepsilon)-h(y;u,\varepsilon)&=u^2(x-y),\\
\mathbb E[|h(x;u,\varepsilon)-\nabla f(x)|^2]
&=\delta^2(x-x^\star)^2+Ls^2,
\end{aligned}
\end{equation}
where $\delta^2=\operatorname{Var}(u^2)$ and $L=\mathbb E[u^2]$.
Label noise increases the variance of fresh gradients but cancels from
differences; feature variability sets $\delta$. Increasing $s$ therefore
affects TSGD's sampled initialization but not its later increments.
When $\delta>0$, the pointwise noise variance is unbounded in $x$,
so similarity holds but global BV fails.
Sections~\ref{subsec:additive-noise-comparison}
and~\ref{subsec:multiplicative-noise-comparison} illustrate the regimes
of exact transport ($\delta=0$) and noiseless labels ($s=0$), respectively.

SGD can also be analyzed under the ABC condition, which generalizes BV
\citep{khaled2023better}. Both constrain the stochastic gradient at
each point separately; centered similarity instead constrains noise
differences at two points under shared sampling.
Appendix~\ref{subsec:comparing-assumptions} compares these conditions.

\subsection{Two Fundamental Sources of TSGD Error}
\label{subsec:martingale-interpretation}

The introduction identified two sources of TSGD error: initialization
and transport. We now make this decomposition precise.
Define $e_k=g_k-\nabla f(x_k)$, $E_k=\mathbb E[\|e_k\|^2]$, and
$d_{k+1}=D_{k+1}-[\nabla f(x_{k+1})-\nabla f(x_k)]$.
The transport update gives
\begin{equation}
\label{eq:estimation-error-recursion}
\boxed{e_{k+1}=e_k+d_{k+1}.}
\end{equation}
Let $\mathcal F_k$ denote the history before drawing $D_{k+1}$,
including $x_{k+1}$ and $g_k$. Conditionally unbiased increments give
\[
\mathbb E[g_{k+1}\mid\mathcal F_k]
=\nabla f(x_{k+1})+e_k.
\]
Unlike SGD, TSGD is therefore generally conditionally biased:
the new increment is unbiased, but the estimator inherits its previous
error. The increment error $d_{k+1}$ has zero conditional mean and is
orthogonal in expectation to previous errors. Consequently,
\begin{equation}
\label{eq:transport-error-identity}
E_k=E_0+\sum_{t=0}^{k-1}\mathbb E[\|d_{t+1}\|^2].
\end{equation}
Even if $\mathbb E[e_0]=0$, so that $\mathbb E[e_k]=0$, the
mean-square error $E_k$ cannot decrease. Initialization error persists,
while transport errors accumulate.
Similarity bounds each contribution by the squared displacement, yielding
\begin{equation}
\label{eq:transport-error-growth}
E_k\leq E_0+\delta^2\sum_{t=0}^{k-1}
\eta_t^2\mathbb E[\|g_t\|^2].
\end{equation}
This explains both the opportunity and the difficulty: accurate changes
preserve useful information, but even small errors can accumulate.
Section~\ref{sec:tsgd-convergence} states the conditional assumptions
directly on $D_{k+1}$, allowing transport beyond pointwise gradient oracles.
